\loai{Cho khối lượng các hạt}
\begin{vidu}%[3P7-4.3K][Loai1]
\nguon{QG - 2017} Trong một phản ứng hạt nhân, tổng khối lượng nghỉ của các hạt trước phản ứng là 37,9638u và tổng khối lượng nghỉ các hạt sau phản ứng là 37,9656u. Lấy $1u=931\ MeV/c^2$. Phản ứng này	
\choice
{tỏa năng lượng 1,68 MeV}
{\True thu năng lượng 1,68 MeV}
{thu năng lượng 16,8 MeV}
{tỏa năng lượng 16,8 MeV}

\loigiai{
Năng lượng phản ứng hạt nhân:
$W=\Delta m.c^2=\left(m{\text{trước}}- m_{Sau}\right)c^2$
\begin{align*}
\Rightarrow W=\left(37,9638 - 37,9656\right)uc^2=\left(37,9638 - 37,9656\right).931,5= - 1,68\ MeV.
\end{align*} 
}
\end{vidu}
\loai{Cho độ hụt khối}
\begin{vidu}%[3P7-4.3B][Loai2]
\nguon{ĐH - 2009} Cho phản ứng hạt nhân: $\Hn{1}{3}{T}+\Hn{1}{2}{D}\rightarrow \Hn{2}{4}{He}+X$. Lấy độ hụt khối của hạt nhân T, hạt nhân D, hạt nhân He lần lượt là 0,009106u; 0,002491u; 0,030382u và $1u = 931,5\ MeV/c^2$. Năng lượng tỏa ra của phản ứng xấp xỉ bằng
\choice
{15,017 MeV}
{200,025 MeV}
{\True 17,498 MeV}
{21,076 MeV}
\loigiai{
Ta có: $W=\left(\Delta m_{\alpha}-\Delta m_T-\Delta m_D\right)c^2=17,498\ MeV$.
}
\end{vidu}
\loai{Cho năng lượng liên kết riêng}
\begin{vidu}%[3P7-4.3K][Loai3]
Cho phản ứng hạt nhân: $\Hn{1}{2}{D}+\Hn{1}{3}{T}\to \Hn{2}{4}{He}+\Hn{0}{1}{n}$. Biết năng lượng liên kết riêng của các hạt nhân tương ứng là: 
${\varepsilon}_{D}=1,1$ $MeV/\text{nuclôn}$, ${\varepsilon}_{T}=2,83$ $MeV/\text{nuclôn}$, ${\varepsilon}_{He}=7,10$ $MeV/\text{nuclôn}$. Năng lượng tỏa ra của phản ứng hạt nhân này là
\choice
{\True 17,71 MeV}
{18,26 MeV}
{17,25 MeV}
{16,52 MeV}	
\loigiai{
Áp dụng công thức: $W=\left(4{\varepsilon}_{He} - 3{\varepsilon}_{T} - 2{\varepsilon}_{D}\right)=17,71\ MeV.$ 
}
\end{vidu}

\loai{Tính năng lượng n mol chất}
\begin{vidu}%[3P7-4.3K][Loai4]
\nguon{QG - 2017} Cho phản ứng hạt nhân: $\Hn{3}{7}{Li}+\Hn{1}{1}{H}\to\Hn{2}{4}{He}+X$. Năng lượng tỏa ra khi tổng hợp được 1 mol helium theo phản ứng này là $5,{2.10}^{24}\ MeV$. Lấy $N_A=6,{023.10}^{23}\ mo{l}^{- 1}$. Năng lượng tỏa ra của một phản ứng hạt nhân trên là	
\choice
{\True 17,3 MeV}
{51,9 MeV}
{34,6 MeV}
{69,2 MeV}
\loigiai{
\begin{itemize}
\item Mỗi phản ứng hạt nhân trên cho ra hai hạt nhân He.
\item Năng lượng tương ứng là: $W = 2\dfrac{{{E_0}}}{{{N_A}}} = 17,3\ MeV.$
\end{itemize}
}
\end{vidu}

\loai{Tính năng lượng tổng hợp m gam chất}
\begin{vidu}%[3P7-4.3K][Loai5]
\nguon{CĐ - 2011} Cho phản ứng hạt nhân $\Hn{1}{2}{H} +\Hn{3}{6}{Li} \rightarrow \Hn{2}{4}{He} +\Hn{2}{4}{He}$. Biết khối lượng các hạt deuterium, liti, helium trong phản ứng trên lần lượt là 2,01360u; 6,01702u; 4,00150u. Coi khối lượng của nguyên tử bằng khối lượng hạt nhân của nó. Lấy $1u=931,5\ MeV/c^2;\ N_A=6,{023.10}^{23}\ mo{l}^{- 1}$. Năng lượng tỏa ra khi có 1 g helium được tạo thành theo phản ứng trên là
\choice
{\True $3,1.10^ {11}\ J$}
{$4,2.10^ {10}\ J$}
{$2,1.10^ {10}\ J$}
{$6,2.10^ {11}\ J$}
\loigiai{
Ta có: $\mathrm {W}_ {\mathrm {\text{tỏa}}} = \dfrac {1} {2} N_ {\alpha} \left(m_ {0} - m \right) c^ {2} = \dfrac {1} {2} \cdot \dfrac {1} {4} \cdot 6,02.10^ {23} \cdot (2,0136 + 6,01702 - 2.4,0015).931,5.1,6.10^ {- 13}=3,1.10^ {11}\ J$.
}
\end{vidu}
\loai{Cho năng lượng toả - thu. Tính năng lượng khi dùng hết m gam chất}
\begin{vidu}%[3P7-4.3K][Loai6]
Cho phản ứng hạt nhân: $\Hn{1}{2}{D}+\Hn{1}{2}{D}\to \Hn{1}{3}{T}+\Hn{1}{1}{H}$. Biết độ hụt khối của các hạt nhân $\Hn{1}{3}{T}$ và $\Hn{1}{2}{D}$ lần lượt là 0,0087u và 0,0024u. Lấy $1u=931\ MeV/c^2;\ N_A=6,{023.10}^{23}\ mo{l}^{- 1}$. Coi khối lượng của nguyên tử bằng khối lượng hạt nhân của nó. Năng lượng tỏa ra trong phản ứng trên khi dùng hết 1 g $\Hn{1}{2}{D}$ là
\choice
{$10,935.10^{23}\ MeV$}
{7,266 MeV}
{\True $5,467.10^{23}\ MeV$}
{3,633 MeV}
\loigiai{
\begin{itemize}
\item Năng lượng tỏa ra trong một phản ứng là: $W=(\Delta {m_T}-2\Delta {m_D})c^2=3,633\ MeV$.
\item Số hạt $\Hn{1}{2}{D}$ có trong 1 g $\Hn{1}{2}{D}$ là: $N=n. {N_A}=\dfrac{m}{A}{N_A}=0,5. 6,02.10^{23}=3,01.10^{23}$ hạt.
\item Một phản ứng cần 2 hạt $\Hn{1}{2}{D}$ tham gia phản ứng, nên năng lượng tỏa ra khi khi dùng hết 1 g $\Hn{1}{2}{D}$ là: $E=W.\dfrac{N}{2}=5,467.10^{23}\ MeV$.
\end{itemize}
}
\end{vidu}
\loai{Bài toán tổng hợp}
\begin{vidu}%[3P7-4.3K][Loai7]
\nguon{QG - 2017} Cho phản ứng hạt nhân $\Hn{6}{12}{C}+\gamma \rightarrow 3\Hn{2}{4}{He}$. Biết khối lượng của $\Hn{6}{12}{C}$ và $\Hn{2}{4}{He}$ lần lượt là 11,9970u và 4,0015u; lấy $1u = 931,5 \ MeV/c^2$. Năng lượng nhỏ nhất của photon ứng với bức xạ $\gamma$ để phản ứng xảy ra có giá trị \textbf{gần nhất} với giá trị nào sau đây?
\choice
{\True 7 MeV}
{6 MeV}
{9 MeV}
{8 MeV}
\loigiai{
Để phản ứng trên có thể xảy ra thì năng lượng bức xạ $\gamma$ tối thiểu là:
\begin{align*}
\varepsilon_{\gamma}=\left(3m_{\alpha}-m_c\right)931,5=6,99\ MeV.
\end{align*}
}
\end{vidu}










